% =============================================================================
%  Student Course Presentation Template
%  Based on project presentation rules (see .claude/rules/presentation-rules.md)
%  Replace all [REMOVE ...] placeholders with your content, then delete the tags.
% =============================================================================
\documentclass[aspectratio=169,11pt]{beamer}
\usetheme{Madrid}
\usecolortheme{default}

% ---------------------------------------------------------------------------
%  Colour palette (professional academic)
% ---------------------------------------------------------------------------
\definecolor{navy}{HTML}{1A365D}        % Headers, titles
\definecolor{darkgray}{HTML}{2D3748}    % Body text
\definecolor{crimson}{HTML}{9B2335}     % Emphasis, key results
\definecolor{forest}{HTML}{276749}      % Positive / adjusted values
\definecolor{steel}{HTML}{4A5568}       % Tertiary, captions
\definecolor{warmgray}{HTML}{E2D8CC}    % Background tint (sparingly)
\definecolor{lightgray}{HTML}{F7FAFC}   % Slide background

% ---------------------------------------------------------------------------
%  Beamer chrome
% ---------------------------------------------------------------------------
\setbeamertemplate{navigation symbols}{}
\setbeamertemplate{footline}[frame number]

\setbeamercolor{frametitle}{fg=navy,bg=white}
\setbeamercolor{title}{fg=navy}
\setbeamercolor{subtitle}{fg=steel}
\setbeamercolor{author}{fg=darkgray}
\setbeamercolor{date}{fg=steel}
\setbeamercolor{normal text}{fg=darkgray}
\setbeamercolor{itemize item}{fg=steel}
\setbeamercolor{itemize subitem}{fg=steel}
\setbeamercolor{alerted text}{fg=crimson}
\setbeamercolor{structure}{fg=navy}

\setbeamerfont{frametitle}{series=\bfseries,size=\large}
\setbeamerfont{title}{series=\bfseries,size=\Large}

% ---------------------------------------------------------------------------
%  Packages
% ---------------------------------------------------------------------------
\usepackage{booktabs}
\usepackage{graphicx}

% ---------------------------------------------------------------------------
%  Custom commands
% ---------------------------------------------------------------------------
% Highlight a key number (crimson, bold)
\newcommand{\emphnum}[1]{{\color{crimson}\textbf{#1}}}

% Positive result number (forest green, bold)
\newcommand{\goodnum}[1]{{\color{forest}\textbf{#1}}}

% Full-slide transition (plain frame, centered text)
\newcommand{\transitionslide}[1]{%
  \begin{frame}[plain]
  \vfill\centering
  {\Large\color{navy}\textbf{#1}}
  \vfill
  \end{frame}
}

% Takeaway box (shaded, centered)
\newcommand{\takeaway}[1]{%
  \begin{center}
  \colorbox{lightgray}{\parbox{0.85\textwidth}{%
    \centering\color{navy}\textbf{#1}%
  }}
  \end{center}
}

% =============================================================================
%  DOCUMENT
% =============================================================================
\begin{document}

% --- Slide 1: Title ---------------------------------------------------------
\begin{frame}[plain]
\vfill
\begin{center}
  {\Large\color{navy}\textbf{[REMOVE Your Presentation Title]}}\\[12pt]
  {\normalsize\color{darkgray}[REMOVE Your Full Name]\\
  Student ID: [REMOVE z1234567]\\[6pt]
  \color{steel}[REMOVE FINS0000] --- [REMOVE Course Name]\\
  [REMOVE Semester, Year]}
\end{center}
\vfill
\end{frame}

% --- Slide 2: Hook -----------------------------------------------------------
% NOTE: Open with one concrete, surprising fact. No bullet lists here.
\begin{frame}{[REMOVE Assertion title: state the surprising fact]}
  \vfill
  \centering
  {\large [REMOVE One sentence or striking statistic that motivates
  the audience to care about your topic.]}
  \vfill
\end{frame}

% --- Slide 3: Research question ----------------------------------------------
% NOTE: State the question clearly and explain why it matters.
\begin{frame}{[REMOVE Why this question matters for \ldots]}
  \begin{itemize}
    \item \textbf{Research question}: [REMOVE State the precise question]
    \item \textbf{Why it matters}: [REMOVE One-sentence stakes]
  \end{itemize}
\end{frame}

% --- Slide 4: Data -----------------------------------------------------------
% NOTE: Title should be an assertion, e.g. "Our sample covers 2,400 stocks".
\begin{frame}{[REMOVE Assertion about your data scope]}
  \begin{itemize}
    \item \textbf{Source}: [REMOVE e.g.\ CRSP, Bloomberg, Yahoo Finance]
    \item \textbf{Sample}: [REMOVE e.g.\ 500 ASX-listed firms]
    \item \textbf{Period}: [REMOVE e.g.\ January 2015 -- December 2025]
    \item \textbf{Frequency}: [REMOVE e.g.\ daily / monthly returns]
  \end{itemize}
\end{frame}

% --- Slide 5: Method ---------------------------------------------------------
\transitionslide{Method}

\begin{frame}{[REMOVE Assertion describing your method]}
  % NOTE: Explain what you did, not textbook theory.
  \begin{enumerate}
    \item [REMOVE Step 1 of your method]
    \item [REMOVE Step 2 of your method]
    \item [REMOVE Step 3 of your method]
  \end{enumerate}
\end{frame}

% --- Slides 7--9: Results (assertion titles) ---------------------------------
\transitionslide{Results}

% Example table slide
\begin{frame}{[REMOVE Assertion title: state the finding]}
  % NOTE: Show only the 2--3 rows that matter. Bold the key coefficient.
  \begin{table}
    \centering\small
    \begin{tabular}{@{}lcc@{}}
      \toprule
      & Return (\%) & $t$-stat \\
      \midrule
      [REMOVE Factor 1] & \emphnum{[REMOVE 0.42]} & [REMOVE 2.31] \\
      [REMOVE Factor 2] & \goodnum{[REMOVE 0.18]}  & [REMOVE 1.96] \\
      \bottomrule
    \end{tabular}
  \end{table}
  \smallskip
  {\small\color{steel} [REMOVE Source note or sample details]}
\end{frame}

% Example figure slide
\begin{frame}{[REMOVE Assertion title: what the chart shows]}
  \begin{center}
    % Replace the placeholder box with your figure:
    %   \includegraphics[width=0.8\textwidth]{figures/your_figure.pdf}
    \fbox{\parbox{0.7\textwidth}{\centering\vspace{3cm}%
      \color{steel}[REMOVE Replace with figure]\vspace{3cm}}}
  \end{center}
  {\small\color{steel} [REMOVE One-sentence caption]}
\end{frame}

% Additional results slide (add or remove as needed)
\begin{frame}{[REMOVE Third assertion about results]}
  \begin{itemize}
    \item [REMOVE Key finding with \emphnum{number}]
    \item [REMOVE Supporting evidence with \goodnum{positive number}]
  \end{itemize}
\end{frame}

% --- Slide 10: Limitations ---------------------------------------------------
\begin{frame}{[REMOVE Strongest objection or key limitations]}
  % NOTE: State the critique fairly, then respond with evidence.
  \begin{itemize}
    \item \textbf{Limitation 1}: [REMOVE e.g.\ survivorship bias]
    \item \textbf{Limitation 2}: [REMOVE e.g.\ short sample period]
    \item \textbf{Robustness}: [REMOVE Results survive X, Y, and Z]
  \end{itemize}
\end{frame}

% --- Slide 11: Takeaway ------------------------------------------------------
% NOTE: End with the ONE thing the audience should remember. No "Questions?" slide.
\begin{frame}{[REMOVE The single takeaway]}
  \vfill
  \takeaway{[REMOVE One sentence the audience remembers after they forget
  everything else.]}
  \vfill
\end{frame}

% --- Backup slides -----------------------------------------------------------
\appendix
\transitionslide{Backup Slides}

\begin{frame}{[REMOVE Backup: additional detail for Q\&A]}
  [REMOVE Content that supports anticipated questions]
\end{frame}

\end{document}
